% Unitat: definició d'espai normat (valencià).

\begin{frame}
\didactatitle{Definició i propietats dels espais normats}

La mètrica usual en $\mathbb{R}$ és la distància euclídea entre dos punts
$x, y \in \mathbb{R}$, donada per
\[
  d(x,y) = |x-y|.
\]
Aquest concepte es generalitza a espais vectorials utilitzant la noció de norma.

\onlynotes{%
  La idea darrere del concepte de norma és generalitzar les propietats més
  útils de la distància euclídea en $\mathbb{R}$: que mesura una cosa no
  negativa, que escala amb el factor d'escala, i que respecta la desigualtat
  triangular.
}

\begin{definition}[Espai normat]
Un \keyterm{espai normat} és un parell $(E, \|\cdot\|)$, on $E$ és un espai
vectorial sobre el cos $\mathbb{K}$ ($\mathbb{R}$ o $\mathbb{C}$) i
$\|\cdot\|$ és una funció anomenada \keyterm{norma} que satisfà:
\begin{itemize}
  \item \textbf{Positivitat:} $\|x\| \geq 0$, i $\|x\| = 0 \iff x = 0$.
  \item \textbf{Homogeneïtat:} $\|\lambda x\| = |\lambda|\,\|x\|$ per a tot
        $\lambda \in \mathbb{K}$ i $x \in E$.
  \item \textbf{Desigualtat triangular:} $\|x+y\| \leq \|x\| + \|y\|$ per a
        tot $x, y \in E$.
\end{itemize}
Quan no hi haja ambigüitat denotarem l'espai normat simplement com $E$.
\end{definition}

\begin{teaching}[Ritme]
Quinze minuts. Val la pena aturar-se en l'homogeneïtat: quasi ningú veu per què
cal el valor absolut en $|\lambda|$ fins que se'ls demana $\lambda = -1$.
\end{teaching}
\end{frame}

\begin{frame}
\didactatitle{Exemples d'espais normats}

Un espai normat permet mesurar la \keyterm{longitud} dels vectors i definir
conceptes com \keyterm{distància} i \keyterm{convergència}.

\dpause

\begin{example}
En $\mathbb{R}^n$ podem definir diferents normes. Les més comunes són
\[
  \|x\|_1 = \sum_{i=1}^n |x_i|,
  \qquad
  \|x\|_2 = \Bigl(\sum_{i=1}^n |x_i|^2\Bigr)^{1/2},
  \qquad
  \|x\|_\infty = \max_{1 \leq i \leq n} |x_i|.
\]

\onlynotes{%
  \begin{itemize}
    \item $\|x\|_2$ és la \keyterm{norma euclídea}, i la seua mètrica associada
          correspon a la distància euclídea.
    \item $\|x\|_1$ s'utilitza en optimització, per la seua simplicitat
          computacional.
    \item $\|x\|_\infty$ és útil en anàlisi numèrica, quan interessa controlar
          el \keyterm{valor màxim} de les components.
  \end{itemize}
}

També podem definir normes en espais més sofisticats, com el de les funcions
contínues en $[a,b]$ amb
\begin{keyformula}
  \|f\| = \max_{x \in [a,b]} |f(x)|.
\end{keyformula}
\end{example}
\end{frame}
